% Chapter 4: System Architecture
\section{Architecture Overview}

The framework consists of several key components working in concert to deliver an interactive learning experience.

\KnowledgePlatformVision

\begin{figure}[H]
    \centering
    \begin{tikzpicture}[
  node distance=1.2cm and 1.4cm,
  core/.style={draw=AccentBlue,fill=AccentBlue!10,rounded corners,minimum width=4.4cm,minimum height=1.05cm,align=center},
  service/.style={draw=AccentTeal,fill=AccentTeal!10,rounded corners,minimum width=3.3cm,minimum height=1.0cm,align=center},
  support/.style={draw=PromptBorder,fill=PromptBackground,rounded corners,minimum width=3.2cm,minimum height=0.95cm,align=center},
  arrow/.style={-{Latex[length=3mm]},thick,draw=AccentBlue}
]
  \node[core] (api) {Knowledge API\\Canonical ontology and concept network};
  \node[support, above left=1.1cm and 1.1cm of api] (standards) {Standards and\\competency mappings};
  \node[support, above right=1.1cm and 1.1cm of api] (planner) {Lesson-plan assembly\\and sequencing service};

  \node[service, below left=1.7cm and 0.2cm of api] (atlas) {\KnowledgeAtlas\\Course-oriented web experience};
  \node[service, below=1.7cm of api] (mobile) {\AlgorithmicaApp\ and \DisciplinaApp\\Gamified mobile pathways};
  \node[service, below right=1.7cm and 0.2cm of api] (landing) {\LandingPageService\\Client information and onboarding};

  \draw[arrow] (standards) -- (api);
  \draw[arrow] (planner) -- (api);
  \draw[arrow] (api) -- (atlas);
  \draw[arrow] (api) -- (mobile);
  \draw[arrow] (api) -- (landing);
  \draw[arrow] (planner) -- (atlas);
  \draw[arrow] (planner) -- (mobile);
\end{tikzpicture}

    \caption{High-level architecture for the shared knowledge API and client-facing services.}
    \label{fig:knowledge-platform-architecture}
\end{figure}

\section{Component Architecture}

\subsection{Ontology Layer}

The ontology layer stores and manages the CS education ontology developed in Chapter~\ref{ch:ontology-design}. A triple store or semantic database maintains RDF representations of knowledge.

\subsection{Reasoning Engine}

The reasoning engine performs semantic queries and inferences over the ontology to generate personalized learning recommendations and pathways.

\subsection{Learning Management Interface}

The interactive interface presents learning content, tracks student progress, and collects feedback. The interface adapts based on ontology-driven reasoning. In the current project vision, \KnowledgeAtlas{} prioritizes course-oriented planning and pathway inspection, while \AlgorithmicaApp{} and \DisciplinaApp{} can emphasize more gamified progression experiences for learners.

\subsection{Student Model}

The student model maintains learner profiles including knowledge state, learning preferences, and performance history. This enables personalization.

\section{Data Flow}

\begin{enumerate}
    \item Student interacts with the learning interface
    \item System updates student model based on interaction
    \item Reasoning engine queries ontology for personalized recommendations
    \item New content or activities are presented
\end{enumerate}

\section{Technology Stack}

\subsection{Backend}

- RDF triple store for ontology storage (e.g., Apache Jena, AllegroGraph)
- SPARQL engine for semantic queries
- RESTful API for frontend communication

\subsection{Frontend}

- Interactive web interface using modern frameworks
- Visualization components for ontology exploration
- Responsive design for accessibility
- A public-facing landing page for informing clients about the service portfolio and value proposition

\section{Integration Considerations}

The architecture supports integration with existing learning management systems and content repositories, ensuring compatibility with institutional infrastructure.
